\section{Exact identities and simultaneous coefficient bounds}\label{app:purchases}
In this appendix $s=u^2$ and $\eta=B+x$. We compare positive multiples
of the four forms
\[
 (f_c\Ps,\ q_c\Qh^2,\ M_e\bc^2,\ r_c\Rh^2).
\]
The fractions $(b_\Ps,b_Q,b_E,b_R)$ bound a row when
the sum of these four weighted forms minus its square is an explicitly
displayed sum of squares.

The elementary identities needed below are
\begin{align}
 \Aop:=(1+U^2)\Ps-\Rh^2
   &=\Bh^2+U^2\Qh^2+U^2(w-UV\Zc)^2,\notag\\
 \E:=(1+V^2)\Ps-V^2\Bh^2
   &=\Qh^2+\Rh^2+V^2(w-UV\Zc)^2 . \label{eq:elementary}
\end{align}
There is also the opposite form
\begin{align}
 \Bop:=2(\Ps+V^2\Qh^2+U^2\Bh^2)-\Rh^2
   &=(E_0+T_0)^2+2(U-V)^2w^2+4UV(w-\Zc)^2\notag\\
   &\quad+2J_0^2+4UV(U-V)^2\Zc^2,\notag\\
 E_0&=\Zc(1-UV),\quad T_0=UV(w-\Zc),\quad
 J_0=T_0-(U-V)^2\Zc. \label{eq:opposite}
\end{align}
All are sums of squares for $U,V\geq0$.

\subsection*{A common bound for the opposite row}
For a nonnegative polynomial multiplier $a_0$, define the fractions
\begin{align}
 \mathfrak b_\Ps(a_0)&=a_0\phi\left[\frac{1+\theta}{f_c}
                                      +\frac{s}{\Sigma yx}\right],&
 \mathfrak b_Q(a_0)&=\frac{2a_0\phi\theta v^2}{q_cy},\notag\\
 \mathfrak b_E(a_0)&=\frac{2a_0\phi s}{\Sigma yx},&
 \mathfrak b_R(a_0)&=\frac{a_0(1-\phi)}{r_c}.
 \label{eq:oppositefractions}
\end{align}
Since $f_c/(\Sigma y)=1-\theta$ and $M_e/\Sigma=\theta$, exactly
\begin{align}
 &f_c\mathfrak b_\Ps(a_0)\Ps+q_c\mathfrak b_Q(a_0)\Qh^2
 +M_e\mathfrak b_E(a_0)\bc^2+r_c\mathfrak b_R(a_0)\Rh^2-a_0\Rh^2
 \notag\\
 &\hspace{30mm}=a_0\phi[(1-\theta)\Aop+\theta\Bop].
 \label{eq:oppositepurchase}
\end{align}
The direct fraction on $\Rh^2$ is weighted by $1-\phi$ before
$r_c^{-1}$ is estimated.

For $O=\sqrt{j}\,Tv(\alpha-c_0s)\Rh$, put
\[
 \bar a=2jT^2v^2(\alpha^2+c_0^2s^2),\qquad c_0=BTCH.
\]
There is the exact additional square
\[
 \bar a\Rh^2-O^2=jT^2v^2(\alpha+c_0s)^2\Rh^2 .
\]
Therefore \eqref{eq:oppositepurchase} with $a_0=\bar a$ bounds $O^2$.
Using \eqref{eq:phi} and \eqref{eq:positivesourcebounds},
its coefficientwise absolute majorants are
\begin{equation}
 \begin{split}
 \mathfrak b_\Ps(\bar a)&\leq_*
             (128/3)Tv^2(1+2s)(\eta^2+s^2),\\
 \mathfrak b_Q(\bar a)&\leq_*32v^4(\eta^2+s^2),\\
 \mathfrak b_E(\bar a)&\leq_*(512/3)Tsv^2(\eta^2+s^2),\\
 \mathfrak b_R(\bar a)&\leq_*(64/3)Tv^2(\eta^2+s^2).
 \end{split}\label{eq:Omajorants}
\end{equation}
Here $\leq_*$ compares absolute coefficients as polynomials in $s,v^2$.

The entire signed term is bounded by taking $a_0=2\kappa$ in
\eqref{eq:oppositepurchase} and adding $2\kappa V^2\Qh^2$.
Indeed
\[
 2\kappa V^2\Qh^2+2\kappa\Rh^2
   -\kappa(V\Qh+\Rh)^2=\kappa(V\Qh-\Rh)^2 .
\]
The four fractions in this bound have majorants
\begin{equation}
 (12T+48Ts,\quad18v^2,\quad96Ts,\quad6T).
 \label{eq:signedmajorants}
\end{equation}
To see this directly, use the four bounds in \eqref{eq:kappabounds}
in \eqref{eq:oppositefractions}; for the $Q$ fraction use
$[4\phi\theta+2]\kappa/(q_cy)\leq18$.

\subsection*{The radial row}
The positive polynomial $\W$ and its permuted version $\W^\pi$ are given
in Appendix~\ref{app:parent}. Define
\[
 \mathcal R_A=\W+4U^2\Aop,\qquad
 \mathcal R_B=2U^2\W^\pi+
       [U\Zc(U^2-1)-(U^2-1)\Bh]^2 .
\]
The expansions in Appendix~\ref{app:parent} show that they are sums of
squares on both state charts.
Let
\begin{align}
 r_\Ps&=(1-\phi)\frac{4x^2y}{f_c}
       +\phi\left[\frac{8\theta xys}{f_c}
                     +\frac{4(x^2+xs+s^2)}{\Sigma}\right],\notag\\
 r_Q&=(1-\phi)\frac{4xys}{q_c}
        +\phi\frac{8\theta ys^2+4(1-\theta)xys}{q_c},\notag\\
 r_E&=\phi\frac{8s^2+2(s-x)^2}{\Sigma},\qquad
 r_R=(1-\phi)\frac{4xys}{r_c}.
 \label{eq:radialfractions}
\end{align}
Then the exact identity is
\begin{align}
 f_cr_\Ps\Ps+q_cr_Q\Qh^2+M_er_E\bc^2+r_cr_R\Rh^2-R_u^2
 =x^2y[(1-\phi)\W+\phi(1-\theta)\mathcal R_A
                                      +\phi\theta\mathcal R_B].
 \label{eq:radialpurchase}
\end{align}
Expansion using \eqref{eq:elementary} and
$Vw(U^2-1)=U\Zc(U^2-1)+(U^2-1)\Bh$ proves the identity.
The following coefficientwise bounds are sufficient:
\begin{equation}
 \begin{split}
 r_\Ps&\leq_*32B+(256/3)s+(512/3)(\eta+s)^2,\\
 r_Q&\leq_*64Bs+(256/3)s^2,\qquad
 r_E\leq_*(1280/3)(\eta+s)^2,\qquad
 r_R\leq_*(128/3)s .
 \end{split}\label{eq:radialmajorants}
\end{equation}
For example $x^2y/f_c\leq8B$, $xy/f_c\leq8/j$ and $y/q_c\leq8/j$.
Apply $\phi/j\leq4/3$ to the terms containing $s$ or $s^2$,
$\phi/\Sigma\leq128/3$ to the remaining terms, and
$(1-\phi)/(1-p)\leq4/3$ to the $r_R$ term.
For $(s-x)^2$, the coefficient majorant is $(s+x)^2$; their difference
$4xs$ is itself a positive polynomial square.

The bound \eqref{eq:Jpurchase} for $16J_\alpha^2$ includes
$32(1-\phi)(\alpha/B)^2R_u^2$.
Thus the effective radial multiplier is the source-only number
\begin{equation}
 L_R=16+32(1-\phi)(\alpha/B)^2
          \leq16+32\eta^2<17 .
 \label{eq:radialmultiplier}
\end{equation}
Its remaining $f_c\Ps$ and $r_c\Rh^2$ fractions are
\begin{equation}
 j_\Ps=\frac{16\phi j^2\alpha^2y(1+s/x)}{f_c},\qquad
 j_R=\frac{32(1-\phi)\alpha^2q_zs}{r_c}.
 \label{eq:Jfractions}
\end{equation}
Their majorants are $16[8B\eta^2+(32/3)\eta^2s]$ and
$16(64/3)\eta^2s$, respectively.

\subsection*{The two remaining helpers}
Put $\gamma=p^2/\Sigma\leq32$. From \eqref{eq:elementary},
\begin{align}
 \gamma M_ev^4\bc^2+\gamma f_c(yv^2+v^4)\Ps-H_1^2
       &=\gamma f_cyv^2\E,\notag\\
 \gamma TM_esv^2\bc^2+\gamma Tf_cs(y+v^2)\Ps-H_2^2
       &=\gamma Tf_cys\E .
 \label{eq:Hhelpers}
\end{align}
Both are polynomial identities, including on the testing axes.
The total fractions for $8H_1^2+8H_2^2$ are therefore
\[
 h_\Ps=8\gamma[yv^2+v^4+Ts(y+v^2)],\qquad
 h_E=8\gamma[v^4+Tsv^2].
\]

\subsection*{All scalar coefficients at once}
There are two equivalent ways to retain the signed square.
The complete-parent decomposition \eqref{eq:coreidentity} has already
used its exact remainder $\kappa(V\Qh-\Rh)^2$, leaving
$2\kappa V^2\Qh^2+2\kappa\Rh^2$ to be bounded.
Define its fractions by \eqref{eq:oppositefractions} with $a_0=2\kappa$,
plus $2\kappa V^2/q_c$ in the $Q$ component. Then the exact total losses are
\begin{equation}
 \begin{pmatrix}b_\Ps\\b_Q\\b_E\\b_R\end{pmatrix}
 =
 L_R\begin{pmatrix}r_\Ps\\r_Q\\r_E\\r_R\end{pmatrix}
 +\begin{pmatrix}j_\Ps\\0\\0\\j_R\end{pmatrix}
 +16\begin{pmatrix}\mathfrak b_\Ps(\bar a)\\
                    \mathfrak b_Q(\bar a)\\
                    \mathfrak b_E(\bar a)\\
                    \mathfrak b_R(\bar a)\end{pmatrix}
 +\begin{pmatrix}\mathfrak b_\Ps(2\kappa)\\
 \mathfrak b_Q(2\kappa)+2\kappa V^2/q_c\\
 \mathfrak b_E(2\kappa)\\\mathfrak b_R(2\kappa)\end{pmatrix}
 +\begin{pmatrix}h_\Ps+g_c/f_c\\0\\h_E\\0\end{pmatrix}.
 \label{eq:totallosses}
\end{equation}
Every replacement in this equality has the exact positive
remainder displayed above; the last $g_c/f_c$ bounds $g_c\Ps$ directly.
The radial term from \eqref{eq:Jpurchase} is included in $L_R$.

For explicit rational upper bounds, set
\[
 S=u_0^2=2^{-24},\qquad V_0=v_0^2=2^{-50},\qquad
 e=2\delta,\qquad D_0=e^2+S^2.
\]
Apply the nonnegative coefficient majorants at $s=S,v^2=V_0$.
The four resulting sums are bounded by
\begin{align}
 A_\Ps={}&17[32\delta+256S/3+512(e+S)^2/3]
 +16[8\delta e^2+32e^2S/3]\notag\\
 &+(2048/3)\delta V_0(1+2S)D_0
 +12\delta+48\delta S \notag\\
 &+256(\delta^2V_0+V_0^2+\delta^3S+\delta SV_0)
 +2(\delta^2+V_0)<9/16,\notag\\
 A_Q={}&17(64\delta S+256S^2/3)+512V_0^2D_0+18V_0<2^{-22},\notag\\
 A_E={}&(21760/3)(e+S)^2+(8192/3)\delta SV_0D_0
       +96\delta S+256(V_0^2+\delta SV_0)<1/16,\notag\\
 A_R={}&(2176/3)S+(1024/3)e^2S
                +(1024/3)\delta V_0D_0+6\delta<1/64 .
 \label{eq:rationaltotals}
\end{align}
These are strict rational comparisons at the displayed powers of two.
The last term in $A_\Ps$ is a conservative bound on $g_c/f_c\leq2y$.
Thus \eqref{eq:rationaltotals} proves \eqref{eq:totalbounds} as coefficient
norm bounds, not only as pointwise inequalities. In particular it proves
positivity of the fixed scalar reserves in the main text.
